\section{The real problem: decided cells}\label{sec:real}

Over $\R$ the mechanisms of \S\ref{sec:first} decide rank one and the
square $m=k$, and duality (Theorem~\ref{thm:duality}) converts corank one
and corank two into rank one and rank two.  Missing are rank two itself and any control on the size.  This section
supplies both by the same device: an operation producing from a design
a smaller design whose dominating subsets transport back.

Deleting a vector of length at most one (\S\ref{sec:light}) lowers the
size by one but moves the remaining vectors, by a congruence; the
transport back is then an identity between quadratic forms, and it is
available only because the deleted vector was short.  Carathéodory
reduction (\S\ref{sec:crystal}) lowers the size to at most $k(k+1)/2+1$
and leaves the vectors where they are; the transport back is immediate,
because $\SC$ in \eqref{eq:dominate} does not involve the weights.  The first operation stalls when every vector is long, and rank two is
where the stalled case can be settled outright
(\S\ref{sec:ranktwo}).

\subsection{Deleting a light vector}\label{sec:light}

\begin{definition}\label{def:light}
A vector of a design is \emph{light} if $\ell_c\le1$ and \emph{heavy} if
$\ell_c>1$.  The design is \emph{all-heavy} if every one of its vectors is
heavy.
\end{definition}

By Lemma~\ref{lem:trace} the $t$-weighted mean of the $\ell_c$ equals $k$,
so at rank $k\ge2$ a design always has a heavy vector; whether it has a
light one is the dichotomy exploited here.  The elementary fact behind the
deletion is the rank-one criterion
\begin{equation}\label{eq:schurone}
  I_k-g^{\vphantom{\mathsf{T}}}g\tp\psd0
  \quad\Longleftrightarrow\quad
  \lVert g\rVert^2\le1 ,
\end{equation}
whose forward direction is the value of the form at $g$, namely
$\lVert g\rVert^2(1-\lVert g\rVert^2)$, and whose converse is
Cauchy--Schwarz: $\langle g,u\rangle^2\le\lVert g\rVert^2\lVert
u\rVert^2\le\lVert u\rVert^2$.

\begin{theorem}[Lean]\label{thm:light}
Let $m\ge2$ and let $\dsn$ be a design at $(m,k)$ possessing a light vector.
If $\Gtz(m-1,k)$ holds, then $\dsn$ admits a dominating $k$-subset.
\end{theorem}

\begin{proof}
Let $\ell_d\le1$.  Since $m\ge2$ and all weights are positive with
$\sum_ct_c=1$, we have $t_d<1$.

\emph{Step 1: the residual Parseval mass is positive definite.}  Put
$\Sigma_d:=I_k-t_d\,g_d^{\vphantom{\mathsf{T}}}g_d\tp$, which by \eqref{eq:parseval} is
$\sum_{c\ne d}t_c\,g_c^{\vphantom{\mathsf{T}}}g_c\tp$.  For $u\ne0$,
\[
  u\tp \Sigma_d u=\lVert u\rVert^2-t_d\langle g_d,u\rangle^2
  \;\ge\;(1-t_d\ell_d)\lVert u\rVert^2\;>\;0 ,
\]
using Cauchy--Schwarz and $t_d\ell_d\le t_d<1$.  Write $\Sigma_d=LL\tp$ with $L$
invertible and set $R:=L^{-\mathsf{T}}$, so that
\begin{equation}\label{eq:whiten}
  R\tp \Sigma_d R=I_k ,\qquad \det R\ne0 .
\end{equation}

\emph{Step 2: the deflated design.}  On the index set
$E:=\{1,\dots,m\}\setminus\{d\}$ put
\[
  g'_c:=\sqrt{1-t_d}\;R\tp g_c ,
  \qquad
  t'_c:=\frac{t_c}{1-t_d}
  \qquad (c\in E).
\]
The weights are positive and sum to one, because
$\sum_{c\in E}t_c=1-t_d$.  For the Parseval condition the factor
$1-t_d$ cancels against the reweighting, and the congruence identity
$(R\tp g)(R\tp g)\tp=R\tp(g^{\vphantom{\mathsf{T}}}g\tp)R$ gives
\[
  \sum_{c\in E}t'_c\,(g'_c)(g'_c)\tp
  \;=\;\sum_{c\in E}t_c\,R\tp g_c^{\vphantom{\mathsf{T}}}g_c\tp R
  \;=\;R\tp \Sigma_d R\;=\;I_k
\]
by \eqref{eq:whiten}.  So $\bigl((g'_c),(t'_c)\bigr)$ is a design at
$(m-1,k)$.

\emph{Step 3: transport.}  By hypothesis some $C\subseteq E$ with
$\lvert C\rvert=k$ dominates in the deflated design.  Its subset sum is
$\sum_{c\in C}(g'_c)(g'_c)\tp=(1-t_d)R\tp\SC R$, so
domination reads
\[
  R\tp\bigl((1-t_d)\SC-\Sigma_d\bigr)R\;\psd\;0 ,
\]
and congruence by the invertible $R$ gives $(1-t_d)\SC-\Sigma_d\psd0$.  The
identity
\begin{equation}\label{eq:reassemble}
  \SC-I_k
  \;=\;\frac{1}{1-t_d}\Bigl[\bigl((1-t_d)\SC-\Sigma_d\bigr)
        \;+\;t_d\bigl(I_k-g_d^{\vphantom{\mathsf{T}}}g_d\tp\bigr)\Bigr]
\end{equation}
is verified by expanding $\Sigma_d$: the two occurrences of
$t_d\,g_d^{\vphantom{\mathsf{T}}}g_d\tp$ cancel and the bracket equals
$(1-t_d)(\SC-I_k)$.  Its first summand is positive semidefinite by the
previous display, its second by \eqref{eq:schurone} and $\ell_d\le1$, and
the scalar is positive.  Hence $C$ dominates in $\dsn$.
\end{proof}

Lightness enters only in the second summand of \eqref{eq:reassemble};
the rest is the congruence and its bookkeeping.  Iterating the theorem removes light vectors one at a time:

\begin{corollary}[Lean]\label{cor:heavy}
Let $k\ge1$.  If every all-heavy design of rank $k$, of whatever size,
admits a dominating $k$-subset, then $\Gtz(m,k)$ holds for every $m$.
\end{corollary}

\begin{proof}
Strong induction on $m$.  A design at $(m,k)$ has $m\ge k$.  If $m=k$,
apply Proposition~\ref{prop:sizerank}.  If $m>k$ and some vector is light,
then $m-1\ge k\ge1$ and $m\ge2$, so Theorem~\ref{thm:light} applies with
the induction hypothesis $\Gtz(m-1,k)$.  Otherwise the design is
all-heavy and the assumption applies.
\end{proof}

The tetrahedron of Example~\ref{ex:tetra} has $\ell_c=3$ for every $c$ and
is therefore all-heavy: the deletion is unavailable precisely at the
design where, by Example~\ref{ex:tetra}, every dominating triple has
$\lmin(\SC)=1$.  This is the first appearance of a pattern which recurs
throughout: the reductions consume the interior of the problem and leave
the extremal locus of \S\ref{sec:ties} standing.

\subsection{Rank two}\label{sec:ranktwo}

Fix a design at $(m,2)$ and write $g_c=(x_c,y_c)$.  Identify $\R^2$ with
$\C$ by $(x,y)\mapsto x+iy$ --- a device for the computation below, with
no relation to the complex problem of \S\ref{sec:complex} --- and let
$\eta_c\in\R^2$ correspond to the square of the complex number attached to
$g_c$:
\begin{equation}\label{eq:doubleangle}
  \eta_c\;:=\;\bigl(x_c^2-y_c^2,\;2x_cy_c\bigr),
  \qquad
  \hat\ell_c\;:=\;\ell_c-2 .
\end{equation}
Squaring doubles arguments and squares moduli, so $\lVert
\eta_c\rVert=\ell_c$; and since the complex number attached to $g_c$ times
the conjugate of the one attached to $g_d$ has real part
$a_{cd}:=\langle g_c,g_d\rangle$ and imaginary part $\pm b_{cd}$, where
\[
  b_{cd}\;:=\;x_cy_d-x_dy_c ,
\]
taking real parts of the square of that product gives
\begin{equation}\label{eq:wdot}
  \langle \eta_c,\eta_d\rangle\;=\;a_{cd}^2-b_{cd}^2 ,
  \qquad
  a_{cd}^2+b_{cd}^2\;=\;\ell_c\ell_d ,
\end{equation}
the second identity being Lagrange's.  The three scalar components of
\eqref{eq:parseval} at $k=2$ read $\sum_ct_cx_c^2=\sum_ct_cy_c^2=1$ and
$\sum_ct_cx_cy_c=0$, so that
\begin{equation}\label{eq:balance2}
  \sum_c t_c\,\eta_c\;=\;0,
  \qquad
  \sum_c t_c\,\ell_c\;=\;2,
  \qquad
  \sum_c t_c\,\hat\ell_c\;=\;0 .
\end{equation}
The design is thus centred in the squared picture, and its centred lengths
are centred as well.  These three identities are the only use made of
\eqref{eq:parseval} in this subsection.

\begin{example}\label{ex:trine2}
Let $u_0,u_1,u_2\in\R^2$ be unit vectors at mutual angles $2\pi/3$ and put
$g_c=\sqrt2\,u_c$, $t_c=\tfrac13$.  From $\sum_cu_c^{\vphantom{\mathsf{T}}}u_c\tp
=\tfrac32I_2$ one gets $\sum_ct_c\,g_c^{\vphantom{\mathsf{T}}}g_c\tp=I_2$, so this
is a design at $(3,2)$, with $\ell_c=2$ and $\hat\ell_c=0$.  It is the member at
$k=2$, $t_c=\tfrac13$ of the extremal family of
Theorem~\ref{thm:corankone}: \eqref{eq:tielaw} reads
$2\cdot\tfrac13\cdot2=\tfrac43=2-1+\tfrac13$.  Each $\eta_c$ has length $2$,
and the three of them again make angles $2\pi/3$; every pair
$C=\{c,d\}$ has
\[
  \SC\;=\;3I_2-2\,u_e^{\vphantom{\mathsf{T}}}u_e\tp,
  \qquad
  \operatorname{spec}\SC=\{1,3\},
\]
where $e$ is the third index.  Every pair dominates, and every pair
dominates with least eigenvalue exactly $1$.  We call this design the rank-two trine.
\end{example}

The remaining argument turns on one scalar invariant of a pair of
indices.

\begin{definition}\label{def:obstruction}
For indices $c,d$ of a design at $(m,2)$ put
\begin{equation}\label{eq:obstruction}
  K_{cd}\;:=\;\langle \eta_c,\eta_d\rangle-\hat\ell_c\hat\ell_d+2 .
\end{equation}
\end{definition}

\begin{lemma}\label{lem:pairobs}
Let $c\ne d$ with $\ell_c+\ell_d>2$.  Then $\{c,d\}$ dominates if and only
if $K_{cd}\le0$.
\end{lemma}

\begin{proof}
Let $M$ be the $2\times2$ matrix with columns $g_c,g_d$, so that
$X:=S_{\{c,d\}}-I_2=MM\tp-I_2$.  For a $2\times2$ matrix $A$ one has
$\det(A-I)=\det A-\tr A+1$; with $A=MM\tp$, $\det A=(\det M)^2=b_{cd}^2$
and $\tr A=\ell_c+\ell_d$, so
\begin{equation}\label{eq:pairdet}
  \det X\;=\;b_{cd}^2-\ell_c-\ell_d+1 .
\end{equation}
On the other hand, eliminating $a_{cd}^2$ from \eqref{eq:wdot} in
\eqref{eq:obstruction} and using $\hat\ell=\ell-2$,
\[
  K_{cd}
  =\ell_c\ell_d-2b_{cd}^2-(\ell_c-2)(\ell_d-2)+2
  =-2\bigl(b_{cd}^2-\ell_c-\ell_d+1\bigr)
  =-2\det X .
\]
Now $\tr X=\ell_c+\ell_d-2>0$.  A symmetric $2\times2$ matrix of positive
trace is positive semidefinite exactly when its determinant is
nonnegative: if $\det X\ge0$ then $X_{00}X_{11}\ge X_{01}^2\ge0$, so
$X_{00}$ and $X_{11}$ have the same sign and the positive trace makes both
nonnegative, whereupon either $X_{00}=0$ --- forcing $X_{01}=0$ and
leaving the diagonal matrix $\diag(0,X_{11})$ --- or
$X_{00}\,x\tp Xx=(X_{00}x_0+X_{01}x_1)^2+(\det X)x_1^2\ge0$; the converse
is the $2\times2$ minor.  Hence $X\psd0$ if and only if $K_{cd}\le0$.
\end{proof}

At the rank-two trine of Example~\ref{ex:trine2} one has $\hat\ell_c=0$ and $\langle
\eta_c,\eta_d\rangle=4\cos(4\pi/3)=-2$ for $c\ne d$, so $K_{cd}=0$ for every
pair: consistent with Lemma~\ref{lem:pairobs} and with
$\det(\SC-I)=0\cdot2=0$.  The inequality $K_{cd}\le0$ is therefore an
equality on a design, and no argument below may improve it to a strict
inequality.

In an all-heavy design every pair has $\ell_c+\ell_d>2$, so if no pair
dominates then Lemma~\ref{lem:pairobs} gives
\begin{equation}\label{eq:allfail}
  K_{cd}>0\qquad\text{for all }c\ne d .
\end{equation}
The contradiction is produced in two steps: a linear identity satisfied
by $K$, and a choice of test vector.

\begin{lemma}\label{lem:rowsum}
For every $c$ one has $\sum_d t_dK_{cd}=2$, and $K_{cd}=K_{dc}$.
Consequently, for every $u\colon\{1,\dots,m\}\to\R$,
\begin{equation}\label{eq:core}
  \sum_{c,d}t_ct_dK_{cd}u_cu_d
  \;=\;2\sum_c t_cu_c^2
      \;-\;\tfrac12\sum_{c,d}t_ct_dK_{cd}(u_c-u_d)^2 .
\end{equation}
If \eqref{eq:allfail} holds and $u$ is not constant, the left side of
\eqref{eq:core} is strictly smaller than $2\sum_ct_cu_c^2$.
\end{lemma}

\begin{proof}
Summing \eqref{eq:obstruction} against $t_d$ and using the three balance
identities \eqref{eq:balance2},
\[
  \sum_d t_dK_{cd}
  =\Bigl\langle \eta_c,\sum_d t_d\eta_d\Bigr\rangle
   -\hat\ell_c\sum_d t_d\hat\ell_d+2\sum_d t_d
  =0-0+2 .
\]
Symmetry of $K$ is symmetry of the inner product.  For \eqref{eq:core},
expand $(u_c-u_d)^2$ and use the row sum twice:
\[
  \sum_{c,d}t_ct_dK_{cd}u_c^2
  =\sum_c t_cu_c^2\sum_d t_dK_{cd}
  =2\sum_c t_cu_c^2 ,
\]
and likewise for $u_d^2$ by symmetry, so the correction term equals
$4\sum_ct_cu_c^2$ minus twice the left side.  The diagonal terms of the
correction vanish; if $u$ is nonconstant there are $c\ne d$ with $u_c\ne
u_d$, and \eqref{eq:allfail} makes the corresponding summand strictly
positive.
\end{proof}

Identity \eqref{eq:core} is the only use of the balance identities
\eqref{eq:balance2}, and it is sharp: the correction term
vanishes identically when $K\equiv0$, as at the rank-two trine.  It remains to
produce a nonconstant $u$ for which the left side of \eqref{eq:core} is at
least $2\sum_ct_cu_c^2$.  The test vectors are the linear functionals of
the squared picture, $u_c=\langle \eta_c,r\rangle$, and the two lemmas below
treat the two possible positions of the family $(\eta_c)$ in the plane.

\begin{lemma}\label{lem:collinear}
Let a design at $(m,2)$ be all-heavy and suppose the vectors $\eta_c$ all lie
on one line through the origin.  Then some pair dominates.
\end{lemma}

\begin{proof}
Let $r\ne0$ be orthogonal to that line and let $r^{\perp}$ be $r$ rotated
by $\pi/2$, so every $\eta_c$ is a multiple of $r^{\perp}$.  Put
$\mu_c:=\langle \eta_c,r^{\perp}\rangle$; then $\lvert\mu_c\rvert=\lVert
\eta_c\rVert\lVert r\rVert=\ell_c\lVert r\rVert\ne0$.  By
\eqref{eq:balance2} $\sum_ct_c\mu_c=0$, and the weights are positive, so
the $\mu_c$ take both signs; choose $\mu_p>0>\mu_q$.  Since
$\eta_p=(\mu_p/\lVert r\rVert^2)\,r^{\perp}$ and likewise for $q$,
\[
  \langle \eta_p,\eta_q\rangle=\frac{\mu_p\mu_q}{\lVert r\rVert^2}<0,
  \qquad
  \lvert\langle \eta_p,\eta_q\rangle\rvert=\lVert \eta_p\rVert\lVert \eta_q\rVert
  =\ell_p\ell_q ,
\]
whence $\langle \eta_p,\eta_q\rangle=-\ell_p\ell_q$ and
\[
  K_{pq}=-\ell_p\ell_q-(\ell_p-2)(\ell_q-2)+2
        =-2(\ell_p-1)(\ell_q-1)<0
\]
because both vectors are heavy.  By Lemma~\ref{lem:pairobs} the pair
$\{p,q\}$ dominates.
\end{proof}

\begin{lemma}\label{lem:spanning}
Let a design at $(m,2)$ be all-heavy and suppose the vectors $\eta_c$ span
$\R^2$.  Then some pair dominates.
\end{lemma}

\begin{proof}
Introduce the second moments of the squared picture,
\[
  \Omega:=\sum_c t_c\,\eta_c^{\vphantom{\mathsf{T}}}\eta_c\tp,
  \qquad
  \zeta:=\sum_c t_c\,\hat\ell_c\eta_c,
  \qquad
  \sigma:=\sum_c t_c\,\hat\ell_c^2 ,
\]
so that $\Omega$ is a symmetric $2\times2$ matrix with $r\tp\Omega
r=\sum_ct_c\langle \eta_c,r\rangle^2$.  This vanishes only if every $\langle
\eta_c,r\rangle$ does, so the spanning hypothesis gives $\Omega\pd0$.  Since
$\lVert \eta_c\rVert=\ell_c$, the second and third identities of
\eqref{eq:balance2} give
\begin{equation}\label{eq:tromega}
  \sigma=\sum_c t_c(\ell_c-2)^2=\sum_c t_c\ell_c^2-4\sum_ct_c\ell_c+4
        =\tr\Omega-4 .
\end{equation}
Put
\[
  \Psi\;:=\;\Omega^2-\zeta\zeta\tp-2\Omega .
\]

\emph{Some direction is nonnegative for $\Psi$.}  Let $\xi:=\Omega^{-1}
\zeta$.  Expanding the square and using
$\sum_ct_c\langle \eta_c,\xi\rangle^2=\xi\tp\Omega\xi=\zeta\tp
\Omega^{-1}\zeta$ and $\sum_ct_c\hat\ell_c\langle \eta_c,\xi\rangle=\langle
\zeta,\xi\rangle=\zeta\tp\Omega^{-1}\zeta$,
\begin{equation}\label{eq:residual}
  \sum_c t_c\bigl(\langle \eta_c,\xi\rangle-\hat\ell_c\bigr)^2
  \;=\;\sigma-\zeta\tp\Omega^{-1}\zeta ,
\end{equation}
so $\zeta\tp\Omega^{-1}\zeta\le\sigma$: the residual of the least-squares
fit of $\hat\ell$ by a linear functional of $\eta$ is nonnegative.  Choose $R$
invertible with $R\tp\Omega R=I_2$, that is $RR\tp=\Omega^{-1}$, and let
$r_1,r_2$ be its columns.  Then $\tr(R\tp\Omega^2R)
=\tr(\Omega^2RR\tp)=\tr\Omega$ and $\tr(R\tp\zeta\zeta\tp R)
=\zeta\tp\Omega^{-1}\zeta$, so by \eqref{eq:tromega} and
\eqref{eq:residual}
\[
  r_1\tp\Psi r_1+r_2\tp\Psi r_2
  =\tr(R\tp\Psi R)
  =\tr\Omega-\zeta\tp\Omega^{-1}\zeta-4
  =\sigma-\zeta\tp\Omega^{-1}\zeta
  \;\ge\;0 .
\]
Hence $r\tp\Psi r\ge0$ for $r:=r_1$ or $r:=r_2$, and $r\ne0$ because $R$
is invertible.

\emph{The test vector.}  Suppose no pair dominates, so
\eqref{eq:allfail} holds, and put $u_c:=\langle \eta_c,r\rangle$.  By
\eqref{eq:balance2}, $\sum_ct_cu_c=0$; were $u$ constant it would be
identically zero, contradicting the spanning hypothesis.  So
Lemma~\ref{lem:rowsum} applies to $u$.  Compute both sides.  First
\[
  \sum_d t_du_dK_{cd}
  =\Bigl\langle \eta_c,\sum_d t_d\langle \eta_d,r\rangle \eta_d\Bigr\rangle
   -\hat\ell_c\sum_d t_d\hat\ell_d\langle \eta_d,r\rangle
   +2\sum_d t_d\langle \eta_d,r\rangle
  =\langle \eta_c,\Omega r\rangle-\hat\ell_c\langle\zeta,r\rangle ,
\]
the last term vanishing by \eqref{eq:balance2}; multiplying by $t_cu_c$
and summing,
\[
  \sum_{c,d}t_ct_dK_{cd}u_cu_d
  =\langle\Omega r,\Omega r\rangle-\langle\zeta,r\rangle^2
  =r\tp\bigl(\Omega^2-\zeta\zeta\tp\bigr)r ,
  \qquad
  \sum_c t_cu_c^2=r\tp\Omega r .
\]
Lemma~\ref{lem:rowsum} now reads $r\tp(\Omega^2-\zeta\zeta\tp)r<2\,
r\tp\Omega r$, that is $r\tp\Psi r<0$, contradicting the choice of $r$.
\end{proof}

\begin{theorem}[Lean]\label{thm:ranktwo}
$\Gtz(m,2)$ holds for every $m$.
\end{theorem}

\begin{proof}
By Corollary~\ref{cor:heavy} it suffices to treat all-heavy designs.  Such
a design has $m\ge2$, and its squared vectors $\eta_c$ either lie on a line
through the origin or span $\R^2$; Lemma~\ref{lem:collinear} and
Lemma~\ref{lem:spanning} produce a dominating pair in the two cases.
\end{proof}

The proof computes no eigenvalue: Lemma~\ref{lem:pairobs} replaces the
least eigenvalue of a $2\times2$ matrix by a determinant, and the rest
is quadratic algebra in the moments of the squared picture.  The
argument is uniform in $m$ and in the weights, as the induction of
Corollary~\ref{cor:heavy} requires, and it is sharp at the rank-two
trine of Example~\ref{ex:trine2}: there $\Omega=2I_2$, $\zeta=0$ and
$\sigma=0$, so $\Psi=\Omega^2-2\Omega=0$, every direction is admissible
in Lemma~\ref{lem:spanning}, and the strict inequality of
Lemma~\ref{lem:rowsum} fails --- as it must, since the trine has
$K\equiv0$.

For uniform weights $t_c=1/m$, Theorem~\ref{thm:ranktwo} is the theorem of
Sengupta and Pautov \cite{SenguptaPautov2026}, whose base case $m=4$ is
due to Nesterenko \cite{Nesterenko2023}; the substitution
\eqref{eq:doubleangle} and the light/heavy split are theirs.  What is
added is the passage to arbitrary weights, which costs the reweighting of
Theorem~\ref{thm:light} in the light branch, and the replacement of the
Perron--Frobenius step of \cite{SenguptaPautov2026} by
Lemma~\ref{lem:rowsum} and the trace computation of
Lemma~\ref{lem:spanning} in the heavy branch.  The designs with $K\equiv0$
at uniform weights are classified in \cite{Nesterenko2026real}: the
squared vectors fall into three clusters of equal vectors, of sizes
$p,q,r$ with $p+q+r=m$, the cluster of size $p$ having length
$\tfrac1{2m}+\tfrac1{2p}$ and similarly for the other two.  The rank-two
trine is the case $p=q=r=1$ at $m=3$: there the common length is
$\tfrac16+\tfrac12=\tfrac23$, which is $\lVert \eta_c\rVert/m=2/3$ for the
vectors of Example~\ref{ex:trine2}, the factor $m^{-1}$ being the passage
from a design at uniform weights to a matrix with orthonormal columns.

\subsection{Corank at most two}\label{sec:corank}

Duality now carries rank two to corank two.

\begin{theorem}[Lean]\label{thm:coranktwo}
$\Gtz(k+2,k)$ holds for every $k\ge2$.
\end{theorem}

\begin{proof}
By Corollary~\ref{cor:dualcell} the statement is equivalent to
$\Gtz(k+2,2)$, which is a case of Theorem~\ref{thm:ranktwo}.
\end{proof}

Unwinding the duality at $r=m-k=2$ makes the reduction explicit.  Let
$\dsn$ be a design at $(k+2,k)$, and put $m:=k+2$.  Its kernel plane
carries an orthonormal basis, arranged as the columns of an $m\times2$
matrix $Z$ with rows $z_c\adj$, and the co-weight rescaling
$w_c=z_c/\sqrt{1-t_c}\in\R^2$ satisfies
\[
  \sum_c(1-t_c)\,w_c^{\vphantom{\mathsf{T}}}w_c\tp\;=\;I_2 .
\]
A $k$-subset is the complement of a pair $\{a,b\}$, and by the chain in
the proof of Theorem~\ref{thm:duality} that $k$-subset dominates in
$\dsn$ if and only if
\begin{equation}\label{eq:pairdom}
  w_a^{\vphantom{\mathsf{T}}}w_a\tp+w_b^{\vphantom{\mathsf{T}}}w_b\tp
  \;\psd\;\Sdual\;=\;\sum_c t_c\,w_c^{\vphantom{\mathsf{T}}}w_c\tp .
\end{equation}
With $R\tp\Sdual R=I_2$ and $h_c:=R\tp w_c$, condition \eqref{eq:pairdom} is
$h_a^{\vphantom{\mathsf{T}}}h_a\tp+h_b^{\vphantom{\mathsf{T}}}h_b\tp\psd I_2$, that is
domination of the pair $\{a,b\}$ in the design $(h_c,t_c)$ at $(k+2,2)$.
So a dominating pair at rank two produces a dominating $k$-subset at
corank two.

The substance of the step is the normalization in \eqref{eq:pairdom}.
The coefficients in \eqref{eq:kernelframe} are the co-weights $1-t_c$,
which sum to $k+1$, whereas the rank-two input requires weights summing to
one.  Rescaling the kernel rows by $t_c$ rather than by $1-t_c$, or
comparing against $I_2$ instead of against $\Sdual$ in \eqref{eq:pairdom},
yields a false statement; Example~\ref{ex:uneven} exhibits the first of
these failures at corank one.

The reduction is field-independent, but its input is not: over $\C$ rank
two fails at size four (Theorem~\ref{thm:complex}), and accordingly
$\GtzC(k+2,k)$ is false for every $k\ge2$ --- the two statements are the
same statement, read through Corollary~\ref{cor:dualcell}.

\begin{corollary}[Lean]\label{cor:corank}
$\Gtz(m,k)$ holds whenever $m-k\le2$, and in particular whenever $m\le5$.
\end{corollary}

\begin{proof}
Corank zero, one and two are Proposition~\ref{prop:sizerank},
Theorem~\ref{thm:corankone} and Theorem~\ref{thm:coranktwo}; the last
requires $k\ge2$, and for $k=1$ every size is covered by
Proposition~\ref{prop:rankone}.  If $m\le5$ then either $k\le2$, whence
Proposition~\ref{prop:rankone} or Theorem~\ref{thm:ranktwo} applies, or
$k\ge3$ and $m\le5\le k+2$ gives $m-k\le2$.
\end{proof}

The two arms of the decided region --- small rank and small corank ---
therefore meet in rank two.  Corank two is also where the description of
the designs attaining the threshold ceases to be a single law.  At corank
one \eqref{eq:tielaw} determines the leverages from the weights; at
$(5,3)$ there are extremal designs whose vectors are pairwise
non-parallel and which satisfy no such relation
(Example~\ref{ex:diamond}).  The corank arm stops here: $(7,4)$ has corank
three and is carried by Corollary~\ref{cor:dualcell} to $(7,3)$, which is
open.

\subsection{Crystallization}\label{sec:crystal}

Write
\[
  \win(k)\;:=\;\frac{k(k+1)}{2}+1 ,
\]
one more than the dimension of the space $\mathrm{Sym}_k$ of real
symmetric $k\times k$ matrices.  Above that size a design carries a
smaller design on a subfamily of its own vectors.

\begin{lemma}[Lean]\label{lem:null}
Let $(g_c)_{c\le m}$ be vectors in $\R^k$ with $m>\win(k)$.  There is
$\delta\in\R^m$, $\delta\ne0$, with
\begin{equation}\label{eq:moment}
  \sum_c \delta_c\,g_c^{\vphantom{\mathsf{T}}}g_c\tp=0
  \qquad\text{and}\qquad
  \sum_c \delta_c=0 .
\end{equation}
\end{lemma}

\begin{proof}
The map $\R^m\to\mathrm{Sym}_k\oplus\R$ sending $\delta$ to
$\bigl(\sum_c\delta_c\,g_c^{\vphantom{\mathsf{T}}}g_c\tp,\sum_c\delta_c\bigr)$ is
linear, and its target has dimension $k(k+1)/2+1=\win(k)<m$.  A linear map
between real vector spaces of dimensions $m$ and $\win(k)<m$ has nonzero
kernel.
\end{proof}

\begin{theorem}[Lean]\label{thm:reduce}
Let $\dsn$ be a design at $(m,k)$ with $m>\win(k)$.  There is a proper subset
$S$ of its index set and a design $\dsn'$ at $(\lvert S\rvert,k)$ whose
vectors are $(g_c)_{c\in S}$, indexed in increasing order.
\end{theorem}

\begin{proof}
Take $\delta$ as in Lemma~\ref{lem:null}.  Since $\sum_c\delta_c=0$ and
$\delta\ne0$, the set $\{c:\delta_c>0\}$ is nonempty; let
\[
  \theta\;:=\;\min\Bigl\{\frac{t_c}{\delta_c}\;:\;\delta_c>0\Bigr\}\;>\;0 ,
\]
attained at $c_0$, and put $t'_c:=t_c-\theta\,\delta_c$.  These numbers are
nonnegative: if $\delta_c\le0$ then $t'_c\ge t_c>0$, and if $\delta_c>0$
then $\theta\le t_c/\delta_c$ gives $\theta\delta_c\le t_c$.  They sum to $1$ and
satisfy $\sum_ct'_c\,g_c^{\vphantom{\mathsf{T}}}g_c\tp=I_k$, by the two identities
\eqref{eq:moment}.  Finally $t'_{c_0}=0$.

Let $S:=\{c:t'_c>0\}$, so $c_0\notin S$ and $S$ is proper.
Restricting the two sums to $S$ omits only vanishing terms, so the
vectors $(g_c)_{c\in S}$ with the weights $(t'_c)_{c\in S}$, re-indexed by
the increasing enumeration of $S$, form a design at $(\lvert S\rvert,k)$.
\end{proof}

\begin{theorem}[Lean]\label{thm:window}
Fix $k$.  If $\Gtz(m,k)$ holds for every $m\le k(k+1)/2+1$, then
$\Gtz(m,k)$ holds for every $m$.
\end{theorem}

\begin{proof}
Strong induction on $m$; the case $m\le \win(k)$ is the hypothesis.  Let
$m>\win(k)$ and let $\dsn$ be a design at $(m,k)$.  Take $S$ and $\dsn'$ as in
Theorem~\ref{thm:reduce} and put $m':=\lvert S\rvert<m$.  By the induction
hypothesis some $C\subseteq S$ with $\lvert C\rvert=k$ dominates in $\dsn'$.
The subset sum
$\SC$ of \eqref{eq:dominate} does not involve the weights and the vectors
of $\dsn'$ are those of $\dsn$, so the subset sum of $C$ in $\dsn$ is the subset
sum of $C$ in $\dsn'$, and $C$ dominates in $\dsn$.
\end{proof}

The bound is Carathéodory's \cite{Caratheodory1911}: condition
\eqref{eq:parseval} exhibits $I_k$ as a convex combination of the points
$g_c^{\vphantom{\mathsf{T}}}g_c\tp$ of the $k(k+1)/2$-dimensional space
$\mathrm{Sym}_k$, and a point of a convex hull in dimension $N$ is a
convex combination of at most $N+1$ of the points (\cite{Rockafellar1970},
\S17).  The proof of Theorem~\ref{thm:reduce} is the standard
support-reduction step of that theorem, carried out by hand for one
reason: what is needed is not some representation of $I_k$ by few
matrices, but a representation by a \emph{subfamily of the given
vectors}, since only then is the dominating subset produced downstairs a
subset of the original index set.  This step uses the weight-freeness of \eqref{eq:dominate}, and it
separates the reduction from that of Theorem~\ref{thm:light}, where the
vectors move and \eqref{eq:reassemble} must carry the conclusion back.

Domination never enters Lemma~\ref{lem:null}, so there is no reason to
expect $\win(k)$ to be the true threshold; see Problem~\ref{prob:higher}.

\subsection{Monotonicity in the size}\label{sec:monotone}

The finitely many cells left by Theorem~\ref{thm:window} are not
independent: the statement at a size implies the statement at all smaller
sizes.  The mechanism is that a design can be padded by splitting one of
its vectors into two equal copies of half the weight, and that a
$k$-subset containing both copies never dominates.

\begin{lemma}[Lean]\label{lem:repeat}
Let $C$ be a subset of the index set of a design at $(m,k)$ with
$\lvert C\rvert\le k$ containing two distinct indices carrying the same
vector.  Then $C$ does not dominate.
\end{lemma}

\begin{proof}
Delete one of the two indices from $C$, obtaining $C_0$ with $\lvert
C_0\rvert\le k-1$.  The linear map $\R^k\to\R^{C_0}$, $p\mapsto(\langle
g_c,p\rangle)_{c\in C_0}$, has nonzero kernel; take $p\ne0$ in it.  Then
$\langle g_c,p\rangle=0$ for every $c\in C$, the deleted index included,
since its vector is one of the retained ones.  Hence
\[
  p\tp(\SC-I_k)p=\sum_{c\in C}\langle g_c,p\rangle^2-\lVert p\rVert^2
  =-\lVert p\rVert^2<0 . \qedhere
\]
\end{proof}

\begin{proposition}[Lean]\label{prop:splitting}
$\Gtz(m+1,k)\Rightarrow\Gtz(m,k)$.
\end{proposition}

\begin{proof}
For $k=0$ there is nothing to prove.  Let $\dsn$ be a design at $(m,k)$, so
$m\ge k$.  Define a design $\dsn^{+}$ at $(m+1,k)$ by keeping the vectors
$g_1,\dots,g_m$ and appending $g_{m+1}:=g_1$, with weights
$t^{+}_1=t^{+}_{m+1}=t_1/2$ and $t^{+}_c=t_c$ for $2\le c\le m$.  Both
$\sum_ct^{+}_c$ and $\sum_ct^{+}_c\,g_c^{\vphantom{\mathsf{T}}}g_c\tp$ are
unchanged, so $\dsn^{+}$ is a design.  By hypothesis some $C^{+}$ with
$\lvert C^{+}\rvert=k$ dominates in $\dsn^{+}$.  By Lemma~\ref{lem:repeat},
$C^{+}$ does not contain both
$1$ and $m+1$, so the map fixing $1,\dots,m$ and sending $m+1$ to $1$ is
injective on $C^{+}$; its image $C$ has $k$ elements and the same subset sum,
hence dominates in $\dsn$.
\end{proof}

\begin{theorem}[Lean]\label{thm:topcell}
$\Gtz\bigl(\win(k),k\bigr)\Rightarrow\Gtz(m,k)$ for every $m$.  Hence the
conjecture at rank $k$ is a single statement, at the single size
$k(k+1)/2+1$.
\end{theorem}

\begin{proof}
Proposition~\ref{prop:splitting}, iterated, gives $\Gtz(m,k)$ for every
$m\le \win(k)$; Theorem~\ref{thm:window} then gives every $m$.
\end{proof}

\subsection{Rank three}\label{sec:rankthree}

At $k=3$ the bound of Theorem~\ref{thm:window} is $M(3)=7$ and
Corollary~\ref{cor:corank} disposes of every size up to five.

\begin{corollary}\label{cor:twocells}
$\Gtz(m,3)$ for all $m$ is equivalent to the conjunction of $\Gtz(6,3)$
and $\Gtz(7,3)$, and also to $\Gtz(7,3)$ alone.
\end{corollary}

\begin{proof}
The first equivalence is Theorem~\ref{thm:window} together with
Corollary~\ref{cor:corank}; the second is Theorem~\ref{thm:topcell} at
$k=3$, or equally Proposition~\ref{prop:splitting} applied once to obtain
$\Gtz(6,3)$ from $\Gtz(7,3)$.
\end{proof}

Both cells are carried through the rest of the paper even though the
second implies the first, because the variational argument of
\S\ref{sec:compactness} at a cell consumes the two cells one size below:
at $(7,3)$ it consumes $(6,3)$ and $(6,2)$.  The logical reduction to one
cell and the analytic route to that cell run in opposite directions.

Neither cell is reducible by the symmetries at hand.  The cell $(6,3)$ is
self-dual, having rank and corank both equal to $3$; and
Theorem~\ref{thm:duality} carries $(7,3)$ to $(7,4)$, again of corank
three.  Since the corank arm stops at two
(Theorem~\ref{thm:coranktwo}), duality cannot shorten the list.

The same pair of cells is reached from above.  Padding a design at $(m,k)$
with one vector along a new coordinate axis produces a design at
$(m+1,k+1)$, and the padding reverses:

\begin{proposition}[Lean]\label{prop:spike}
Let $1\le k\le m$.  Then $\Gtz(m+1,k+1)\Rightarrow\Gtz(m,k)$.
\end{proposition}

\begin{proof}
Let $\dsn$ be a design at $(m,k)$ in which no $k$-subset dominates.  For each
of the finitely many $k$-subsets $C$ there is, by \eqref{eq:dominate}, a
vector $p_C\ne0$ with $\sum_{c\in C}\langle g_c,p_C\rangle^2<\lVert
p_C\rVert^2$; let $\sigma\in(0,1)$ be the average of $1$ and the largest
of the finitely many ratios $\sum_{c\in C}\langle
g_c,p_C\rangle^2/\lVert p_C\rVert^2$, so that
\begin{equation}\label{eq:unifgap}
  \sum_{c\in C}\langle g_c,p_C\rangle^2<\sigma\lVert p_C\rVert^2
  \qquad\text{for every $k$-subset }C .
\end{equation}
Build a design at $(m+1,k+1)$: the vectors
$\hat g_c:=(\sigma^{-1/2}g_c,0)$ with weights $\sigma t_c$ for $c\le m$,
together with $\hat g_{m+1}:=(0,\dots,0,(1-\sigma)^{-1/2})$ of weight
$1-\sigma$.  The weights are positive and sum to one, and
\eqref{eq:parseval} holds blockwise, the two scalings cancelling in each
block.  By hypothesis some $C^{+}$ with $\lvert C^{+}\rvert=k+1$ dominates there.
If $m+1\notin C^{+}$,
every selected vector has last coordinate zero and the last standard basis
vector gives $1\le0$.  So $m+1\in C^{+}$ and $C:=C^{+}\setminus\{m+1\}$ has $k$
elements; testing domination against $(p_C,0)$ gives
\[
  \lVert p_C\rVert^2\;\le\;\sigma^{-1}\sum_{c\in C}\langle
  g_c,p_C\rangle^2\;<\;\lVert p_C\rVert^2
\]
by \eqref{eq:unifgap}, a contradiction.  Hence some $k$-subset of $\dsn$
dominates after all.
\end{proof}

Combining Proposition~\ref{prop:spike} with Corollary~\ref{cor:twocells},
$\Gtz(8,4)$ implies all of rank three.  The family of statements
$\Gtz(m,k)$ is therefore monotone in both parameters, and no implication
in this section can be shown to be strict: were the conjecture true, all
of them would be equivalences, so a strictness proof would refute it.

\begin{figure}[ht]
  \centering
  \begin{tikzpicture}[
  x=1cm, y=1cm,
  every node/.style={font=\small, inner sep=3pt},
  proved/.style={draw=black!55, line width=0.4pt, fill=black!10,
                 rounded corners=1pt, align=center},
  openq/.style ={draw=black, line width=0.9pt, fill=white,
                 rounded corners=1pt, align=center},
  hyp/.style   ={font=\scriptsize, align=center, inner sep=1pt},
  tag/.style   ={font=\scriptsize, inner sep=1.5pt},
  imp/.style   ={-{Latex[length=2.1mm,width=1.5mm]}, line width=0.5pt},
  bnd/.style   ={-{Triangle[open,length=2.2mm,width=1.9mm]}, line width=0.5pt},
  cnj/.style   ={-{Latex[length=2.1mm,width=1.5mm]}, line width=0.5pt,
                 dash pattern=on 1.4pt off 1.2pt}]

\node[proved] (two) at (2.7,2.05) {rank two: $\Gtz(m,2)$ for every $m$};
\node[proved] (low) at (0.9,0)    {corank ${\le}\,2$\\[1pt] $(3,3)\ \ (4,3)\ \ (5,3)$};
\node[openq]  (six) at (5.0,0)    {$(6,3)$};
\node[openq]  (sev) at (8.3,0)    {$(7,3)$};
\node[openq]  (all) at (11.4,0)   {$\Gtz(m,3)$\\[1pt] for every $m$};

\node[hyp] (h6) at (5.0,-1.75) {interior\\ exclusion};
\node[hyp] (h7) at (8.3,-1.75) {interior\\ exclusion};

% the boundary analysis at a cell consumes the two cells one size below
\draw[bnd] (low.east) -- node[tag, above] {$(5,3)$} (six.west);
\draw[bnd] (two.south) to[out=-90,in=120]
     node[tag, pos=0.30, left=2pt] {$(5,2)$} (six.north);
\draw[bnd] (two.east) to[out=0,in=105]
     node[tag, pos=0.55, above right=-1pt] {$(6,2)$} (sev.north);
\draw[bnd] ([yshift=2.5pt]six.east) --
     node[tag, above] {boundary} ([yshift=2.5pt]sev.west);

% the logical reductions
\draw[imp] ([yshift=-2.5pt]sev.west) --
     node[tag, below] {splitting} ([yshift=-2.5pt]six.east);
\draw[imp] (sev.east) -- node[tag, above] {$M(3)=7$} (all.west);

\draw[cnj] (h6) -- (six);
\draw[cnj] (h7) -- (sev);

% ---------------------------- legend --------------------------------
\begin{scope}[shift={(0,-3.05)}]
  \draw[imp] (0,0) -- (0.9,0);
  \node[tag, anchor=west] at (1.05,0) {implies};
  \draw[bnd] (3.1,0) -- (4.0,0);
  \node[tag, anchor=west] at (4.15,0) {supplies the boundary case};
  \draw[cnj] (8.5,0) -- (9.4,0);
  \node[tag, anchor=west] at (9.55,0) {conjectural input};
\end{scope}

\end{tikzpicture}

  \caption{Rank three.  Solid arrows are the reductions of
  Corollary~\ref{cor:corank} and Theorem~\ref{thm:window}; the boundary
  analysis of \S\ref{sec:compactness} at a cell consumes the two cells one
  size below.  Dashed arrows carry the two unproved inputs of
  \S\ref{sec:interior}.}
  \label{fig:ladder}
\end{figure}
